% =====================================================================
%  CHAPTER 7 -- CONCLUSION AND FUTURE WORK
%
%  COMMENT / DEPARTMENT VARIATION:
%  Some departments want a single "Conclusion" chapter and treat future
%  work as a closing paragraph rather than a separate section; others
%  ask for a distinct "Future Prospects" section. This chapter is
%  written with the two split into Sections 7.1 and 7.2. If your
%  department does not require a separate future-work section, merge
%  Section 7.2 into Section 7.1 and delete the \section line -- nothing
%  else in the project depends on it.
%
%  RUBRIC NOTE (R9): the top band wants a succinct summary of ALL
%  essential results, summarised both quantitatively and qualitatively,
%  an objective account of strengths and limitations, and a clear and
%  complete set of recommendations for follow-up work.
% =====================================================================
\chapter{CONCLUSION AND FUTURE WORK}
\label{ch:conclusion}

% ---------------------------------------------------------------------
\section{CONCLUSION}
\label{sec:conclusion}
% ---------------------------------------------------------------------

% TODO: restate the problem and what was built, in two or three
% sentences. No new information belongs in a conclusion.
[RESTATEMENT: recall the problem stated in
Section~\ref{sec:problem-statement} and state what was built to address
it.]

% TODO: summarise the results QUANTITATIVELY. Repeat the key numbers
% from Chapter 6 -- the rubric asks for exactly this.
[QUANTITATIVE SUMMARY: restate the headline measured results. Give the
actual numbers again here; do not send the reader back to Chapter 6.]

% TODO: summarise QUALITATIVELY -- what the work demonstrates.
[QUALITATIVE SUMMARY: what the project demonstrates in broader terms,
and what its strengths are.]

% TODO: reflect briefly and honestly on the design process itself.
[REFLECTION: a short, objective reflection on the design process -- what
worked, what would be done differently, and what the main limitation of
the delivered system is.]

% ---------------------------------------------------------------------
\section{FUTURE WORK}
\label{sec:future-work}
% ---------------------------------------------------------------------

% TODO: give CONCRETE recommendations. "Improve the UI" is not a
% recommendation; "add offline synchronisation so field users can record
% inspections without connectivity" is.
Building on the work presented in this report, the following
directions are recommended for future development:

% NOTE: the braces around each bracketed placeholder are REQUIRED --
% without them LaTeX reads "[...]" as a custom label for \item.
\begin{enumerate}[label=\arabic*., leftmargin=*, itemsep=2pt]
  \item {[RECOMMENDATION 1: a concrete functional extension, and why it
        would be valuable.]}
  \item {[RECOMMENDATION 2: a technical improvement -- scalability,
        performance, security hardening -- and what it would require.]}
  \item {[RECOMMENDATION 3: an evaluation the present work could not
        carry out, such as a larger field trial.]}
  \item {[RECOMMENDATION 4: add or remove entries as appropriate.]}
\end{enumerate}
